\documentclass[12pt, a4paper]{article}

% ---------- Encoding & Language ----------
\usepackage[utf8]{inputenc}
\usepackage[T1]{fontenc}
\usepackage[ngerman]{babel}

% ---------- Fonts ----------
\usepackage{lmodern}

% ---------- Layout ----------
\usepackage[
  top=2.5cm,
  bottom=2.5cm,
  left=3cm,
  right=2.5cm
]{geometry}
\setlength{\parskip}{6pt}
\setlength{\parindent}{0pt}

% ---------- Header / Footer ----------
\usepackage{fancyhdr}
\setlength{\headheight}{15pt}
\addtolength{\topmargin}{-3pt}
\pagestyle{fancy}
\fancyhf{}
\fancyhead[L]{\small\leftmark}
\fancyhead[R]{\small\thepage}
\renewcommand{\headrulewidth}{0.4pt}

% ---------- Sections styling ----------
\usepackage{titlesec}

% Section: groß, fett, mit Linie darunter
\titleformat{\section}
  {\large\bfseries}
  {\thesection}{1em}{}
  [\vspace{2pt}\hrule height 0.8pt\vspace{4pt}]
\titlespacing{\section}{0pt}{14pt}{6pt}

% Subsection: normal fett, deutlich sichtbar
\titleformat{\subsection}
  {\normalsize\bfseries}
  {\thesubsection}{1em}{}
  [\vspace{1pt}\hrule height 0.3pt\vspace{2pt}]
\titlespacing{\subsection}{0pt}{10pt}{4pt}

% Subsubsection: normalsize, fett
\titleformat{\subsubsection}
  {\normalsize\bfseries}
  {\thesubsubsection}{1em}{}
\titlespacing{\subsubsection}{0pt}{8pt}{2pt}

% ---------- Boxes ----------
\usepackage[most]{tcolorbox}

\newtcolorbox{definitionbox}[1][]{
  colback=white, colframe=black, coltext=black,
  fonttitle=\bfseries, title={Definition}, #1
}
\newtcolorbox{notebox}[1][]{
  colback=white, colframe=black, coltext=black,
  fonttitle=\bfseries, title={Hinweis}, #1
}
\newtcolorbox{examplebox}[1][]{
  colback=white, colframe=black, coltext=black,
  fonttitle=\bfseries, title={Beispiel}, #1
}

% ---------- TikZ ----------
\usepackage{tikz}
\usetikzlibrary{arrows.meta, positioning, shapes.geometric, fit, backgrounds, decorations.pathreplacing}

% ---------- Math ----------
\usepackage{amsmath, amssymb}

% ---------- Lists ----------
\usepackage{enumitem}
\setlist[itemize]{leftmargin=1.5em, itemsep=2pt}
\setlist[enumerate]{leftmargin=1.5em, itemsep=2pt}

% ---------- Figures ----------
\usepackage{float}

% ---------- Hyperlinks ----------
\usepackage[hidelinks]{hyperref}

% ---------- Misc ----------
\usepackage{microtype}

% =========================================
%                DOCUMENT
% =========================================
\begin{document}

\begin{center}
  {\LARGE\bfseries Präsentationstechnik}\\[6pt]
  {\large Vorlesungsnotizen}\\[4pt]
  {\small Stand: \today}
\end{center}

\vspace{1em}
\tableofcontents
\newpage

% =========================================
\section{Grundlagen der Kommunikation}
% =========================================

Der erste Block klärt, was beim Sprechen zwischen zwei Menschen überhaupt
passiert. Die Modelle sind das Handwerkszeug für alles Weitere: wer weiß,
warum eine Aussage anders ankommt als sie gemeint war, kann seinen Vortrag
darauf einstellen.

\subsection{Sender-Empfänger-Modell}

Das Modell stammt von Shannon und Weaver. Der Sender \emph{denkt},
\emph{codiert} und \emph{sendet}; der Empfänger \emph{empfängt},
\emph{decodiert} und \emph{versteht}. Zurück läuft eine Antwort, so dass ein
Kreis entsteht.

\begin{itemize}
  \item Beide Seiten haben einen \textbf{begrenzten Zeichenvorrat}. Was der
        eine nicht ausdrücken und der andere nicht auflösen kann, geht
        verloren.
  \item Codieren und Decodieren hängen von \textbf{Persönlichkeit und
        Sozialisation} ab. Dasselbe Wort trägt bei zwei Menschen nicht
        dieselbe Bedeutung.
  \item Auf dem Weg wirken \textbf{Störungen}.
  \item Übergeordnete Rahmenbedingungen: zeitlich, sozial, physikalisch
        (Hitze, Gerüche, Licht) und kulturell.
\end{itemize}

\begin{notebox}
Das Modell erklärt, warum Missverständnisse der Normalfall und nicht die
Ausnahme sind: die Nachricht wird zweimal übersetzt, und beide Übersetzungen
benutzen verschiedene Wörterbücher.
\end{notebox}

\subsection{Die fünf Axiome von Paul Watzlawick}

\begin{enumerate}
  \item \textbf{Man kann nicht \emph{nicht} kommunizieren.} Jedes Verhalten
        ist Kommunikation; sie beginnt, sobald Menschen einander wahrnehmen.
        Schweigen und Wegsehen sind ebenfalls Botschaften.
  \item \textbf{Jede Kommunikation hat einen Inhalts- und einen
        Beziehungsaspekt.} Der Inhaltsaspekt ist die WAS-Botschaft, reine
        Sachinformation, überwiegend verbal übermittelt. Der
        Beziehungsaspekt ist die WIE-Botschaft, verbal und nonverbal, und er
        gibt vor, wie die Sachinformation aufzufassen ist. Widersprechen
        sich beide direkt, wird die Kommunikation schwierig.
  \item \textbf{Kommunikation ist immer Ursache und Wirkung.} Sie ist
        zirkulär; was Ursache und was Wirkung ist, liegt in der subjektiven
        Interpretation der Beteiligten. Grundlage ist der Konstruktivismus:
        wir leben in einer konstruierten Wirklichkeit, geformt aus
        persönlichen Erfahrungen und Werten, und diese bestimmt unser
        Handeln. Das Problem entsteht, wenn zwei Interaktionspartner von
        verschiedenen Wirklichkeiten ausgehen.
  \item \textbf{Menschliche Kommunikation bedient sich analoger und digitaler
        Modalitäten.} Digital ist das Miteinander-Sprechen: es verfügt über
        logische Operatoren, es mangelt ihm aber an Wortschatz für die
        Beziehungsebene. Analog sind Körpersprache, Sprechweise und Kontext:
        reich an Beziehungsinformation, dafür nicht eindeutig. Beide
        ergänzen sich und sind getrennt voneinander schwer zu interpretieren;
        eine Diskrepanz stört die Kommunikation.
  \item \textbf{Kommunikation ist symmetrisch oder komplementär.}
        Symmetrisch heißt gleichwertig, die Beteiligten versuchen
        Ungleichheit zu vermeiden. Komplementär heißt ergänzend,
        unterschiedliche Verhaltensweisen ergänzen sich und bestimmen so den
        Interaktionsprozess. Wir leben in vielen Rollen und wechseln
        zwischen beidem.
\end{enumerate}

\subsection{Die vier Seiten einer Nachricht (Schulz von Thun)}

Jede Äußerung hat vier Seiten. Der Sender hat vier Schnäbel, der Empfänger
vier Ohren, und ob \emph{gemeint} gleich \emph{aufgenommen} ist, entscheidet
sich daran, welches Ohr gerade groß ist.

\begin{itemize}
  \item \textbf{Sachinformation:} worüber ich informiere.
  \item \textbf{Beziehungsbotschaft:} was ich von dir halte und wie wir
        zueinander stehen.
  \item \textbf{Selbstkundgabe:} was ich von mir zu erkennen gebe.
  \item \textbf{Appell:} wozu ich dich veranlassen möchte.
\end{itemize}

\begin{examplebox}
\textbf{\glqq Was ist das Grüne in der Soße?\grqq}

\emph{Er} meint:
Sachebene \glqq Da ist was Grünes\grqq{},
Selbstoffenbarung \glqq Ich weiß nicht, was es ist\grqq{},
Beziehung \glqq Du wirst es wissen\grqq{},
Appell \glqq Sag mir, was es ist!\grqq{}

\emph{Sie} versteht:
Sachebene \glqq Da ist was Grünes\grqq{},
Selbstoffenbarung \glqq Mir schmeckt das nicht\grqq{},
Beziehung \glqq Du bist eine miese Köchin!\grqq{},
Appell \glqq Lass nächstes Mal das Grüne weg!\grqq{}

Die Antwort \glqq Wenn es dir hier nicht schmeckt, kannst du ja woanders
essen gehen!\grqq{} ist die Reaktion auf die empfangene, nicht auf die
gesendete Nachricht.
\end{examplebox}

\subsection{Transaktionsanalyse von Eric Berne}

Ziel: Konflikte im Gespräch lösen, das Gesprächsverhalten interpretieren und
steuern, das Gesprächsklima verbessern. Kleinste Einheit ist die
\textbf{Transaktion} aus Transaktions-Stimulus und Transaktions-Reaktion.

\subsubsection{Die Ich-Zustände}

\begin{itemize}
  \item \textbf{kEL} Kontrollierendes Eltern-Ich: tadelnd, anweisend,
        bewertend, kritisierend, moralisierend, verallgemeinernd.
  \item \textbf{fEL} Fürsorgliches Eltern-Ich: lobend, umsorgend,
        beschützend, bevormundend, behütend, schonend.
  \item \textbf{ER} Erwachsenen-Ich: zuhörend, fragend, feststellend,
        überprüfend, kritisch, emotionslos, sachlich.
  \item \textbf{rK} Rebellisches Kind-Ich: aufbegehrend, schert sich nicht um
        Konventionen.
  \item \textbf{aK} Angepasstes Kind-Ich: freundlich, brav, zuverlässig,
        pflichtbewusst.
  \item \textbf{fK} Freies Kind-Ich: bedürfnisorientiert,
        begeisterungsfähig, kontaktfreudig, Konventionen spielen keine
        Rolle.
\end{itemize}

\subsubsection{Die drei Transaktionsarten}

\begin{itemize}
  \item \textbf{Parallele (einfache, komplementäre) Transaktion:} der Partner
        reagiert aus dem Ich-Zustand, aus dem der andere gesprochen hat. Das
        Gespräch läuft glatt.
  \item \textbf{Gekreuzte Transaktion:} die Reaktion kommt nicht aus dem
        angesprochenen Ich-Zustand und ist unerwartet. Hier entstehen
        Konflikte.
  \item \textbf{Verdeckte Transaktion:} die Gesprächspartner sagen etwas
        anderes, als sie meinen. Neben der offenen läuft eine zweite,
        verdeckte Botschaft.
\end{itemize}

\subsection{Konstruktives Feedback}

Feedback ist die individuelle und \textbf{subjektive Reaktion} auf einen
Menschen: \glqq So nehme ich Sie wahr, und so reagiere ich darauf.\grqq{} Es
hat drei Elemente: \textbf{Wahrnehmung}, \textbf{Interpretation},
\textbf{Reaktion}.

\subsubsection{Feedback geben}

\begin{itemize}
  \item Negatives Feedback von positivem flankieren.
  \item Konkret und präzise bleiben, keine Verallgemeinerungen oder
        pauschalen Aussagen.
  \item Auf professioneller Ebene kommunizieren, Kritik sachlich äußern.
  \item Subjektiv formulieren: \glqq Ich-Botschaften\grqq{}.
  \item Konstruktiv sein, also Perspektiven aufzeigen, aber nicht
        herablassend.
  \item Raum für Rückmeldungen geben und diese annehmen: offene
        Formulierungen, Dialog auf Augenhöhe.
\end{itemize}

\subsubsection{Feedback annehmen}

\begin{itemize}
  \item Ausreden lassen und aktiv zuhören.
  \item Wiederholen, was man verstanden hat.
  \item Bei Unklarheiten nachfragen, aber die Aussagen nicht in Frage
        stellen.
  \item \textbf{Nicht rechtfertigen.}
  \item Sich bedanken.
  \item Das Feedback akzeptieren und in Ruhe, mit Abstand reflektieren: eine
        Chance zur kritischen Selbstreflexion.
\end{itemize}

\subsection{Stress und Lampenfieber}

Öffentlich zu reden ist mit \textbf{41\,\%} die häufigste Angst der Menschen,
vor großen Höhen (32\,\%), Geldmangel, tiefem Wasser und Ungeziefer (je
22\,\%). Subjektiv wahrgenommen wird sie als \textbf{Leistungsangst} und als
\textbf{soziale Angst}. Eine Präsentation beansprucht die Ressourcen stark,
die Folge ist Stress.

\subsubsection{Stresslevel und Leistung}

Die Kurve ist ein umgekehrtes U. Bereich A ist zu wenig Erregung: wenig
Energie, Unterspannung, im Extrem Schlaf. Bereich B ist das Optimum mit hoher
Aufmerksamkeit. Bereich C ist Überregung: Kontrollverlust, Angst und Panik.

\begin{notebox}
Ein wenig Lampenfieber ist keine Störung, sondern der Arbeitspunkt. Ziel ist
nicht null Erregung, sondern die Mitte der Kurve.
\end{notebox}

\subsubsection{Drei Zugänge zum Stressbegriff}

\begin{itemize}
  \item \textbf{Stress als Input}, Situationsansatz von Holmes und Rahe
        (1967): Stress hängt von externen Umweltbedingungen ab, den
        \emph{Stressoren}, etwa kritischen Lebensereignissen. Diese lassen
        sich nach ihrem stressbedingten Krankheitswert vergleichen. Da jeder
        Mensch anders auf Umweltreize reagiert, gibt es keinen
        allgemeingültigen Katalog von Stressfaktoren.
  \item \textbf{Stress als Transaktion}, Verarbeitungsprozesse nach Lazarus
        (1999): 1. primäre Bewertung, die Stressoren werden bewertet;
        2. sekundäre Bewertung, die eigenen Bewältigungsmöglichkeiten werden
        eingeschätzt; 3. beide Ergebnisse werden zusammengeführt. Personen-
        und Umweltmerkmale spielen eine Rolle, also Nature (genetische
        Anlagen) und Nurture (Erfahrung).
  \item \textbf{Stress als Reaktion}, Selye (1936): Stress ist die
        unspezifische Reaktion des Organismus auf Anforderungen aus der
        Umwelt. Allgemeines Adaptionssyndrom in drei Phasen: 1. Alarmphase,
        2. Widerstandsphase, 3. krankhafter Erschöpfungszustand.
\end{itemize}

Zusammengesetzt: Stressoren (Input) $\rightarrow$ Verarbeitung (Transaktion)
$\rightarrow$ Reaktionen (Output).

\subsubsection{Umgang mit Stress}

Ressourcenaktivierung (was habe ich in der Vergangenheit bereits geschafft,
und wie?), positive Glaubenssätze und kognitive Umstrukturierung, Planung und
Zeitmanagement, eine Notfallstrategie, Entspannung, Sport und Genießen
lernen. Über allem steht die soziale Unterstützung.

% =========================================
\section{Vortragsplanung}
% =========================================

\subsection{Präsentationssituation: die Checkliste vorab}

\begin{itemize}
  \item \textbf{Vorbereitung:} Rahmen der Präsentation, klare
        Zielvorstellungen (eigene Ziele, Ziele des Auftraggebers,
        Erwartungen des Publikums), Stoff sammeln.
  \item \textbf{Inhalt:} Strukturierung und Gliederung, Notizen.
  \item \textbf{Vortragsart:} üben und Feedback holen, Körpersprache
        einsetzen.
  \item \textbf{Organisation:} Raum, Zeitpunkt.
  \item \textbf{Technik:} Formate klären, mehrere Medien, Back-up der
        Präsentation.
  \item \textbf{Für die Präsentation:} Handout, Back-up-Folien für die
        Diskussion, ein Glas Wasser.
\end{itemize}

\subsection{Rahmen und Redeart}

Zu klären ist, in welchem Rahmen die Präsentation stattfindet, wie formell
oder informell er ist und welchen Ton der Vortrag haben soll:
Informationsrede, Argumentationsrede (begeistern, überzeugen) oder
Gelegenheitsrede zu einem bestimmten Anlass.

\subsection{Drei Zielgruppen}

Die Ziele kommen aus drei Richtungen: den \textbf{eigenen Zielen}, den Zielen
des \textbf{Auftraggebers} und den Erwartungen des \textbf{Publikums}.

\begin{examplebox}
Beim Referat: der \emph{Dozent} will Studierende bewerten können, die eigene
Veranstaltung strukturieren und das Gesagte für die Diskussion nutzen. Der
\emph{Student} will die Prüfung ableisten, eine gute Note bekommen, das
Referat hinter sich bringen und sich nicht blamieren. Die \emph{Kommilitonen}
wollen etwas für sich mitnehmen, unterhalten werden oder die Zeit absitzen.
\end{examplebox}

\subsection{Ziele}

\begin{itemize}
  \item \textbf{Ziele bestimmen die Inhalte.} Sie bilden den roten Faden und
        halten den Vortrag zusammen.
  \item Ein Ziel sagt, was \emph{erreicht} werden soll, nicht, was ich sagen
        möchte.
  \item Gibt es eine \textbf{Zielhierarchie}? Was sind zentrale Ziele, was
        Nebenziele?
  \item \textbf{Anzahl der Ziele:} wie viele lassen sich in der Zeit
        erreichen? Zu viele Ziele überfrachten den Vortrag.
\end{itemize}

\begin{definitionbox}
\textbf{SMARTe Ziele}\\
\textbf{S}pezifisch: konkret, unmissverständlich, detailliert.\\
\textbf{M}essbar: qualitativ und quantitativ.\\
\textbf{A}ttraktiv: angemessen, erreichbar, ausführbar.\\
\textbf{R}ealistisch: finanziell, technisch, zeitlich.\\
\textbf{T}erminiert: Zeitrahmen, was bis wann?
\end{definitionbox}

\subsection{Erwartungen der Zuhörer}

Zu fragen ist: Wer ist das Publikum, wer nimmt teil und wie viele? Was ist
seine Motivation, warum sind die Zuhörer dort, ist die Veranstaltung
freiwillig? Wie steht das Publikum zum Thema, welche Haltungen gibt es, was
weiß es und was nicht?

Der Zuhörer darf verlangen, dass Vortragende ihr Thema beherrschen, den
Zusammenhang zwischen Aussagen und Zielen verdeutlichen, Aussagen angemessen
und verständlich darstellen, die Zuhörer respektieren, den Vortrag auf deren
Kenntnisstand abstimmen und Aufmerksamkeit und Interesse aufrecht halten.

\subsection{Inhalte planen in vier Phasen}

\begin{enumerate}
  \item \textbf{Auflisten:} Wie lassen sich die Ziele erreichen? Relevante
        Aussagen sammeln, Kreativitätstechniken wie Mind Mapping einsetzen.
  \item \textbf{Sortieren:} die gesammelten Argumente, Beispiele und Ideen
        strukturieren. Was sind zentrale Punkte, wie lassen sie sich
        veranschaulichen? Was fehlt, was kann weg? Zeit und zentrale Aussage
        im Auge behalten.
  \item \textbf{Gliedern:} Ablauf festlegen, Kernbotschaften herausarbeiten,
        den Vortrag nicht überfrachten (sonst gehen Prägnanz und
        Verständlichkeit verloren), Einleitungs- und Schlussteil entwerfen.
  \item \textbf{Überarbeiten:} Sind die Übergänge logisch? Ist der zeitliche
        Rahmen realistisch, oder müsste ich durchhetzen? Sprechernotizen als
        Stichpunkte und Regieanweisungen anlegen. Wenn möglich ganz frei
        vortragen.
\end{enumerate}

\subsection{Aufbau und Zeitaufwand}

\begin{center}
\begin{tabular}{lr}
\textbf{Teil} & \textbf{Anteil}\\\hline
Eröffnung & 15\,\%\\
Hauptteil & 75\,\%\\
Schluss & 10\,\%\\
\end{tabular}
\end{center}

Weitere Gliederungsmuster: chronologisch, räumlich, kausal, vergleichend,
thematisch, lösungsorientiert.

\subsection{Orientierung}

Der Vortrag muss klar strukturiert sein, und die Struktur muss während des
gesamten Vortrags sichtbar bleiben. Es muss klar sein, warum welcher Punkt
wann kommt und wo der Vortragende gerade steht.

\begin{itemize}
  \item Überblick am Anfang.
  \item Neue Redeabschnitte deutlich machen.
  \item Wegweiser nutzen: \glqq Zunächst möchte ich Ihnen A zeigen, um dann
        zu verdeutlichen, wie dies mit Punkt B zusammenhängt\grqq{},
        \glqq Hier ein Beispiel\grqq{}, \glqq An dieser Stelle möchte ich
        etwas hinzufügen.\grqq{}
\end{itemize}

\subsection{Die Einleitung}

Sie klärt Organisatorisches und die Dauer, weckt Aufmerksamkeit und Interesse
und baut eine Beziehung zu den Zuhörern auf. Sie beantwortet: Worum geht es
und was ist das Ziel? Warum ist das Thema wichtig? Was soll man mitnehmen?
Wie ist der Vortrag gegliedert?

\subsubsection{Einstiegsarten}

\begin{itemize}
  \item \textbf{Orientierende Vortragseröffnung:} Überblick geben, Personal
        Story.
  \item \textbf{Problem- und Nutzeneinstiege:} Knobel-Denkanreiz,
        \glqq Nutzenlocker\grqq{}.
  \item \textbf{Der aktuelle Bezug:} brandaktuell, Anknüpfen an die
        Geschichte der Gruppe.
\end{itemize}

\subsubsection{Aufmerksamkeit wecken: DOs}

Mit einer Frage beginnen; eine Aufgabe oder Handlungsaufforderung geben; eine
gemeinsame Erinnerung oder Erfahrung ansprechen; eine passende, kurze
Anekdote oder ein Zitat; Neuigkeiten; Anschauungsobjekte.

\begin{notebox}
Die ersten Worte bestimmen maßgeblich die Haltung der Zuhörer. Ein guter
Einstieg macht neugierig und gespannt auf mehr.
\end{notebox}

\subsubsection{DON'Ts beim Einstieg}

Sich entschuldigen; Witze oder Geschichten, die nichts mit dem Thema zu tun
haben; eine banale, ideenlos vorgetragene Feststellung; eine langatmige
Aussage oder Geschichte; abgedroschene Phrasen.

\subsection{Der Hauptteil}

Anfang und Schluss lassen sich leichter gestalten, wenn der Hauptteil schon
steht. Hier gilt: nicht überfrachten. Die Kernbotschaften ausdrücken, mit
Argumenten untermauern und mit Beispielen veranschaulichen.

\begin{center}
\begin{tabular}{ll}
\textbf{Empirische Arbeit} & \textbf{Literaturarbeit}\\\hline
Problemstellung & Problemstellung\\
Zielsetzung & Zielsetzung\\
Theorie & Theorie I\\
Vorgehensweise & Theorie II\\
Empirische Ergebnisse & Theorie III\\
Schlussfolgerung \& Empfehlungen & Zusammenführung \& Schlussfolgerung\\
\end{tabular}
\end{center}

\subsection{Beziehung zum Publikum}

Enthusiasmus verbreiten, Augenkontakt halten, einbindende Formulierungen
nutzen (\glqq Wir sollten in Betracht ziehen\ldots\grqq{}), das Publikum
einbeziehen, Feedback suchen.

Verloren geht der Bezug durch: zu viel und unnatürlich lächeln; das Publikum
nicht ansehen; sich immer wieder die Folien ansehen und dem Publikum den
Rücken zeigen; keine Fragen zulassen; eine monotone Stimme; den Vortrag
vorlesen; genervt oder verängstigt auf Fragen reagieren; sich auf Kosten
dummer Fragen profilieren; sich für Fehler rechtfertigen.

\subsection{Der Abschluss}

Summen bilden und den Vortrag abschließen, Perspektiven aufzeigen, im
Gedächtnis bleiben. Bei wissenschaftlichen Vorträgen folgen Nachfragen,
Kritik und die Abschlussdiskussion.

\begin{itemize}
  \item \textbf{Schlussaussage:} Was ist das Fazit, das die Zuhörer mitnehmen
        können?
  \item \textbf{Schlusssatz:} ähnliche Stilmittel wie in der Einleitung, also
        Frage, Zitat oder Handlungsaufforderung.
\end{itemize}

\begin{notebox}
Das Ende bleibt den Zuhörern am längsten in Erinnerung. Hier ist die Chance,
die Kernbotschaften zu verstärken.
\end{notebox}

Zu vermeiden: den Zeitrahmen überziehen; auf das Ende zu hasten (vorher
klären, was bei Zeitnot wegfallen kann); das Ende immer wieder antäuschen und
doch weitermachen; eine weitschweifende Zusammenfassung mit viel
Wiederholung; mit \glqq Ich danke Ihnen für Ihre Aufmerksamkeit\grqq{} enden,
weil das dem Vortrag die Wirkung nimmt. Besser: danken und dann einen
aussagekräftigen Schlusssatz.

% =========================================
\section{Präsentation}
% =========================================

\begin{definitionbox}
Eine \textbf{Präsentation} ist der mündliche Vortrag einer oder mehrerer
Personen, bei dem einem Publikum bestimmte Inhalte in strukturierter Form
unter Verwendung visueller Hilfsmittel dargeboten werden.
\end{definitionbox}

\subsection{Arten, Medien und Programme}

\textbf{Präsentationsarten:} Informationspräsentation; Überzeugungs-,
Meinungs- und Argumentationspräsentation; Anlass- und Gelegenheitsrede.

\textbf{Präsentationsmedien} unterstützen das gesprochene Wort, niemals
umgekehrt. Neben dem Beamer: Overhead-Projektor, Flipchart, Pinnwand,
Whiteboard oder Tafel, Plakat.

\textbf{Programme:} PowerPoint, Prezi, prezentit, SlideRocket, Keynote,
Google Slides, Empressr.

\subsection{Folienanzahl}

So viele wie nötig, so wenige wie möglich. Bei guter Visualisierung sind mehr
Folien möglich als bei reinen Textfolien. Maximal sieben Folien am Stück,
dann eine \glqq Folienpause\grqq{}.

\subsection{Visualisierung}

\subsubsection{Warum}

Visualisieren heißt \glqq etwas bildhaft darstellen\grqq{}. Es erhöht die
Verständlichkeit erheblich, Präsentationen leben von Bild und Text, und die
meisten Menschen sind visuelle Typen.

\textbf{Ziele:} die Aufmerksamkeit auf das Wesentliche konzentrieren; den
Betrachter einbeziehen; Orientierungshilfen geben; Informationen leichter
erfassbar machen; Wesentliches verdeutlichen; Gesagtes ergänzen und
vertiefen; das Behalten fördern.

\begin{center}
\begin{tabular}{lr}
\textbf{Sinne} & \textbf{Behaltensquote}\\\hline
Hören & 20\,\%\\
Hören und Sehen & 50\,\%\\
Selbst tun & 100\,\%\\
\end{tabular}
\end{center}

\subsubsection{Fünf Grundsätze für visuelle Hilfsmittel}

\begin{enumerate}
  \item Der logische Zusammenhang muss erkennbar bleiben.
  \item Beim Umschalten Zeit lassen und neue Informationskanäle ankündigen.
  \item Weniger Visualisierung ist zumeist mehr.
  \item Jedes Bild muss eine bestimmte Funktion haben.
  \item Folien sind kein Ersatz für ein Stichwortkonzept.
\end{enumerate}

\subsubsection{Planung einer Visualisierung}

Erst die drei Fragen: Was will ich darstellen (Inhalt)? Wozu soll die
Darstellung dienen (Ziel)? Wen will ich informieren oder überzeugen
(Zielgruppe)? Erst dann die Visualisierung im engeren Sinne, also wie und
womit.

Bausteine: inhaltliche Elemente $+$ Präsentationsmedien $+$
Gestaltungselemente. Die Gestaltungselemente sind Text, Grafik und Symbole
sowie Diagramme.

\subsection{Text auf der Folie}

Auf gute Lesbarkeit achten: einfache Schriftarten, Lesegewohnheiten
beachten.

\begin{definitionbox}
\textbf{Die vier Verständlichmacher}
\begin{enumerate}
  \item Einfachheit
  \item Gliederung und Ordnung
  \item Kürze und Prägnanz
  \item Zusätzliche Stimulanzen (Farben, Beispiele, Grafiken)
\end{enumerate}
\end{definitionbox}

\subsubsection{Schriftart}

Maximal zwei Schriftarten. Serifenfreie Schriften (etwa Arial), weil
Serifenschriften (etwa Times) projiziert schlechter lesbar sind. Nur wenige,
bedachte Schriftwechsel. Nicht ausschließlich Großbuchstaben, keine
senkrechte Anordnung der Buchstaben.

\subsubsection{Schriftgröße}

Die Schriftgröße wird dem Raum angepasst: \textbf{16\,pt für 5\,m Distanz},
\textbf{24\,pt für 10\,m Distanz}.

\begin{examplebox}
\textbf{Rechenbeispiel.} Monitor: 0,5\,m Abstand, 1024$\times$768 Pixel,
10\,pt $=$ 0,353\,cm. Projektion: 5\,m Abstand, Bild 240$\times$180\,cm. Die
Schrift wird um den Faktor 10 vergrößert, also 3,53\,cm; umgerechnet
768\,px $\cdot$ 3{,}53\,cm $/$ 180\,cm $=$ 15\,px, also 15\,pt. Eine
15-pt-Schrift aus 5\,m wirkt damit wie eine 10-pt-Schrift aus 0,5\,m.
\end{examplebox}

\subsection{Grafik, Symbole, Animation und Medien}

Grafiken verdeutlichen und heben Informationen hervor und lockern trockene
Themen auf, sollen aber nicht \glqq zu viel des Guten\grqq{} sein und den
Inhalt nicht verdrängen.

Bei Animationen, Effekten und Übergängen sparsam vorgehen. Am einfachsten und
oft am besten: Erscheinen und Verschwinden. Zu viele Effekte lenken vom
Vortragenden ab, das Publikum blickt auf die Folie und die Aufmerksamkeit
gilt der Animation; das wirkt schnell unprofessionell. Eingebettete Videos
und Audios sollen professionell und nicht zu lang sein und zum Ziel
beitragen.

\subsection{Diagramme}

\begin{definitionbox}
\textbf{Diagramme} sind standardisierte Darstellungsformen für bestimmte
Sachverhalte.
\end{definitionbox}

Am besten selbst erstellen und auf Screenshots verzichten. Überschrift und
gegebenenfalls Quellenangabe gehören immer dazu.

\begin{itemize}
  \item \textbf{Liste und Tabelle:} machen Zahlen und Werte transparent und
        verschaffen einen Überblick. Bedarfsorientiert, das Wichtigste
        hervorheben (etwa durch Umrahmung), Spaltenüberschriften sollen sich
        entsprechen und optisch abheben, bei vielen Zeilen nummerieren,
        nicht zu sehr ins Detail gehen.
  \item \textbf{Kurvendiagramm:} zeigt Entwicklungsverläufe und vergleicht
        Prozesse. Ausgangspunkt ist der Nullpunkt des Koordinatensystems,
        die horizontale Achse (Abszisse) trägt in der Regel den zeitlichen
        Ablauf. Die Achseneinteilung bestimmt den Kurvenverlauf, jede Achse
        wird klar bezeichnet, bei mehreren Kurven verschiedene Linienarten.
  \item \textbf{Säulen- und Balkendiagramm:} vergleichende Darstellung von
        Größen. Entscheiden, ob Absolutzahlen, Prozentwerte oder kumulierte
        Werte die Aussage am deutlichsten machen; Achseneinteilung sorgfältig
        wählen; gleiche Breite; Abstand höchstens die Säulenbreite;
        Nulllinien deutlich beschriften und die Skalierung angeben.
  \item \textbf{Kreis- oder Tortendiagramm:} stellt das Ganze und seine Teile
        dar. 360 Grad $=$ 100\,\%. Keine zu kleinen Teilmengen, mehrere
        kleine zu einer \glqq Sammeleinheit\grqq{} zusammenfassen, Teile
        optisch klar trennen.
  \item \textbf{Organigramm bzw.\ Aufbaudiagramm:} zeigt Aufgabenverteilung
        und Hierarchien mittels Kästchen; Zusammenhänge durch Lage,
        Strichstärken und Schraffuren. Klassisch von oben nach unten
        detaillierter, alternativ vom Mittelpunkt ausgehend.
  \item \textbf{Ablaufdiagramm bzw.\ Netzplan:} zeigt vorbestimmte Soll- oder
        ermittelte Ist-Abläufe mit einfachen, normierten Symbolen.
\end{itemize}

\subsubsection{Entscheidungs-Matrix}

\begin{center}
\begin{tabular}{ll}
\textbf{Darzustellende Information} & \textbf{Diagrammart}\\\hline
Aufzählung & Liste\\
Datenzuordnung & Tabelle\\
Absolute Werte & Liste, Säulen, Balken\\
Anteile eines Ganzen & Kreis\\
Organisationsstrukturen & Organigramm, Aufbau\\
Aufbau, Zusammensetzung & Organigramm, Aufbau\\
Entwicklungsverläufe, Abläufe & Kurven, Ablauf\\
Gegenüberstellung & Säulen, Balken\\
\end{tabular}
\end{center}

\subsection{Komposition und Farben}

Die Gestaltungselemente werden zu einer präsentationsreifen Vorlage
zusammengestellt. Drei Punkte:

\begin{itemize}
  \item \textbf{Aufteilung:} für eine klare Grobstruktur, hilfreich ist ein
        Raster. Muster sind Reihung, Symmetrie, Rhythmus und Dynamik.
  \item \textbf{Anordnung und Logik:} nicht zufällig, sondern die logische
        Struktur dessen widerspiegelnd, was dargestellt werden soll.
  \item \textbf{Farben und Formen} als Bedeutungsträger: wichtige
        Informationen hervorheben, Zusammenhänge verdeutlichen, Querverweise
        herstellen, aufeinanderfolgende Darstellungen verbinden.
        Kontrastreiche, gut lesbare Farben; \textbf{kein Gelb oder Grün für
        Text}; sparsam und bewusst, \textbf{höchstens drei Farben}; möglichst
        in einer Farbpalette bleiben.
\end{itemize}

\subsubsection{Hintergrund}

\begin{center}
\begin{tabular}{p{0.45\textwidth}p{0.45\textwidth}}
\textbf{Weiß auf Schwarz} & \textbf{Schwarz auf Weiß}\\\hline
auffällig; gut für Bildschirmpräsentationen; Lesbarkeit hängt vom Medium ab;
nicht gut bei den meisten Beamern und in hellen Räumen
&
leichte Lesbarkeit, auch bei schwächeren Beamern und in helleren Räumen;
Handouts lassen sich leichter ausdrucken\\
\end{tabular}
\end{center}

Grafiken im Hintergrund dürfen nicht ablenken, sonst lieber ganz darauf
verzichten.

\subsection{Struktur der Folie und allgemeine Tipps}

Nur 5 bis 7 Punkte pro Folie (situationsabhängig: Studentenvortrag gegen
klausurrelevanten Dozentenvortrag). Nicht zu viel Text, Symbole und Pfeile.
Schnell erfassbare Überschriften. Nur die wesentlichen Aussagen in
Stichworten. Foliennummer und Namen nicht vergessen.

Allgemein: Blöcke bilden; gleiche Farben und Formen für zusammengehörende
Sachverhalte; Wichtiges hervorheben; Freiflächen als Gestaltungselement;
\glqq im Weglassen liegt die Kunst\grqq{}.

Weitere Tipps: Zeit investieren; Folien ankündigen und Überleitungen
schaffen; ein Bild sagt mehr als 1000 Worte; mehr sagen als das, was auf den
Folien steht; PowerPoint lernt sich nur in der praktischen Übung.

\subsection{Checkliste}

Gliederung und Ablauf; kurze, knappe Sätze oder Stichworte; Schriftart (Arial,
Verdana); Schriftgröße; optische Sinnabschnitte; einheitliche Folien;
Masterfolie benutzt; Gliederung auf jeder Folie; Grafik und Bilder dienen dem
Inhalt; nicht mehr als drei Farben; hoher Kontrast zwischen Hintergrund und
Text; konsistente Seitenzahlen; unter realen Bedingungen durchgeklickt.

% =========================================
\section{Rhetorik}
% =========================================

\begin{definitionbox}
\textbf{Rhetorik} ist Theorie und Praxis mündlicher Kommunikation. Angewandte
Rhetorik ist im engeren Sinne die Rede- und Gesprächspädagogik. Ihre
Geburtsstätte ist Griechenland; im 5.\ Jahrhundert v.\ Chr.\ entstand die
eigentliche Rhetorik.
\end{definitionbox}

Fachliches Know-how (Hard Skills) und rhetorisches Geschick (Soft Skills)
gehören zusammen; das eine trägt ohne das andere nicht.

\subsection{Warum Rhetorik wichtig ist}

\begin{enumerate}
  \item Wir sind auf das gemeinsame Gespräch angewiesen.
  \item Was wir sagen, ist oft nicht identisch mit dem, was wir sagen
        wollten.
  \item Was der Zuhörer aufnimmt, ist nicht immer identisch mit dem wirklich
        Gesagten.
  \item Richtiges setzt sich nicht allein deshalb durch, weil es richtig ist;
        es muss anderen auch als richtig erklärt werden.
\end{enumerate}

Das eigene Sprechen verbessert man in drei Schritten: rhetorische Prozesse
erkennen lernen, die Techniken der Kommunikation kennen, den Transfer in die
Praxis üben.

\subsection{Was Rhetorik umfasst}

\begin{center}
\begin{tabular}{lll}
\multicolumn{2}{c}{\textbf{Nonverbale Kommunikation}} &
\textbf{Verbale Kommunikation}\\\hline
\textbf{Körpersprache} & \textbf{Räumliches Verhalten} & \textbf{Tonalität}\\
Körperhaltung, Gestik, & Distanzverhalten, & Stimme und Stimmklang,\\
Blickkontakt, Mimik & Körperorientierung, & Aussprache, Betonung\\
 & Bewegung im Raum & \\
\end{tabular}
\end{center}

\subsection{Die drei Säulen der Rhetorik}

\begin{itemize}
  \item \textbf{Ethos}, Vertrauen durch persönliche Werte: Reputation,
        Expertise, Zitate, Publikumsinteraktion, Gemeinsamkeiten.
  \item \textbf{Pathos}, die Welt der Emotionen: Vision und Träume,
        Metaphern, Humor, Geschichten, Verletzlichkeit.
  \item \textbf{Logos}, Fakten: Fakten und Daten, Forschung und Statistiken,
        rhetorische Fragen, Demonstrationen, Testergebnisse.
\end{itemize}

\subsection{Körperhaltung}

\textbf{Der Stand:} mit beiden Füßen fest auf dem Boden stehen, das Gewicht
gleichmäßig verteilen, gelegentlich Stand- und Spielbein wechseln. Das ist
ein Zeichen für eine gute Mischung aus Standhaftigkeit und Flexibilität. Die
Bewertungsskala reicht von \emph{zurückgezogen} über \emph{unsicher} bis
\emph{selbstbewusst}.

\textbf{Das Sitzen:} Oberkörper möglichst aufrecht halten, Schultern locker,
aber nicht hängen lassen, keine Barriere mit den Beinen bilden, Raum
einnehmen und selbstbewusst wirken, jedoch nicht zu viel. Die Skala reicht
von \emph{ablehnend} über \emph{verkrampft} bis \emph{souverän}.

\subsubsection{Die vier Distanzen}

\begin{center}
\begin{tabular}{ll}
Intime Distanz & $<$ 60\,cm\\
Persönliche Distanz & 60--120\,cm\\
Gesellschaftliche Distanz & 1--4\,m\\
Öffentliche Distanz & $>$ 4\,m\\
\end{tabular}
\end{center}

Für einen Vortrag vor einer Gruppe mindestens die persönliche Distanz wählen.
Achtung: Distanzen sind kulturabhängig.

\subsection{Gestik}

Arme, Hände und Finger nutzen; punktuell bestimmte Aussagen unterstreichen;
Demonstrationsgesten nutzen, indem man Gesagtes bildhaft darstellt;
Aufzählungen mit den Fingern mitzählen. Auch Gesten können kulturabhängig
sein.

Die Körpersprache soll offen sein, also keine verschränkten Arme. Hände in
Bauchhöhe strahlen konzentrierte Ruhe aus. Zu vermeiden sind nervöses
Trommeln mit den Fingern und verkrampfte, verknotete Finger; möglichst oft
die Handflächen zeigen. Hand in der Hosentasche, Stift oder Manuskript in der
Hand: auch hier gilt es, die Balance zu finden. Gesten lassen den Vortrag
lebendiger wirken, zu viel Gestik lässt ihn hektisch und unruhig wirken. Die
eigene Position findet man durch Üben vor dem Spiegel.

\subsection{Mimik}

Menschen unterscheiden ein echtes und ein aufgesetztes Lächeln in der Regel
sehr gut. Der Trick besteht darin, an etwas zu denken, das einen lächeln
lässt. Eine freundliche, offene und interessierte Mimik entsteht aus einer
freundlichen, offenen und interessierten Haltung.

\textbf{Funktionen der Mimik:} sie spiegelt persönliche Eigenschaften wider,
zeigt die innere Stellungnahme des Einzelnen, zeigt Emotionen und regelt
Interaktionsabläufe.

\textbf{Sechs Basisemotionen} (kulturunabhängig): Freude, Wut, Furcht, Ekel,
Traurigkeit, Überraschung. Ziel ist die Kongruenz zwischen dem Gesagten und
der Mimik, denn nur so wirkt der Vortrag glaubhaft.

\textbf{Blickkontakt:} gerader Kopf mit offenem Blick, Augenkontakt suchen,
aber nicht starren; sich freundliche oder bekannte Gesichter suchen.

\subsection{Tonalität}

\begin{itemize}
  \item Klar und deutlich sprechen.
  \item Die Stimme modulieren, eine Sprachmelodie einsetzen.
  \item Laut, aber nicht schreien; permanent lautes Sprechen wirkt schnell
        aggressiv.
  \item Zu leise wirkt schüchtern und strengt an, zu hoch wirkt angespannt.
        Tiefere Stimmlagen werden oft als angenehmer empfunden.
\end{itemize}

\textbf{Langsam sprechen:} 90\,\% aller Menschen sprechen zu schnell. Beim
Vortragen muss das normale Sprechtempo reduziert werden, sonst werden
undeutlich ausgesprochene Wörter verschluckt, und die Zuhörer haben keine
Chance, Verpasstes nachzuholen.

\textbf{Richtig atmen:} Sprechen ist Ausatmen. Anspannung führt zu flacher
Atmung und damit zu gepresster Stimme, gehetztem Sprechen, Angst und
Schwindel. Bewusstes Atmen dient dem besseren Sprechen, einer festeren
Stimme und dem Abbau von Stress.

\textbf{Sprechstil:} mit Bildern arbeiten (Beispiele, Veranschaulichung,
Metaphern, Erfahrungen) und positive Formulierungen nutzen
(\glqq Ich freue mich\ldots\grqq{}, \glqq Chance und Gelegenheit\grqq{}).

\textbf{Pausen} sind genauso wichtig wie das Reden. Vier Arten:
Nachwirkungspause, Vorbereitungspause, Regenerationspause, Denkpause.

\subsection{Zu vermeidende Signale}

\begin{center}
\begin{tabular}{ll}
\textbf{Signal} & \textbf{Wirkung}\\\hline
zugekniffene Augen & Abwehrhaltung\\
seltener Blickkontakt & Unsicherheit, Verlegenheit\\
häufiger Lidschlag, Lidflattern & Unsicherheit\\
angehobene Augenbrauen & Ungläubigkeit, Arroganz\\
geöffneter Mund & mangelnde Selbstkontrolle\\
verschränkte Arme & Ablehnung, Angst\\
aneinander reibende Hände & Selbstgefälligkeit\\
zu einem Spitzdach geformte Hände & Arroganz, Abwehr\\
geballte Fäuste & Wut, Aggression\\
häufiges Achselzucken & Hilflosigkeit, Ungewissheit\\
wippende Füße & Ängstlichkeit, Ungeduld\\
\end{tabular}
\end{center}

\subsection{Die fünf Rednertypen}

Rhetorik muss zum Redner, seiner Persönlichkeit und seinen Fähigkeiten
passen. Man kann mehreren Typen zugleich zuzuordnen sein, und Mischformen
sind die beste Voraussetzung für einen gelungenen Vortrag.

\begin{itemize}
  \item \textbf{Der mitteilungsfreudige Rednertyp.} Ausprägungen: der
        Kommunikative sucht den Dialog zum Publikum, der rein
        Mitteilungsfreudige verfällt schnell in einen Monolog. Strahlt
        Offenheit aus, schweift leicht ab und verliert Anlass und Ziel aus
        den Augen. Tipps: genauer Plan über Inhalte und verfügbare Zeit, den
        Vortrag in Kontrollabschnitte unterteilen.
  \item \textbf{Der schüchterne Rednertyp.} Leidet unter Lampenfieber und
        Selbstzweifel und muss die Nervosität in den Griff bekommen. Tipps:
        sicherer Einstieg, indem die ersten Zeilen 1:1 auswendig gelernt
        werden; vorher mit Zuschauern plaudern; freundliche und interessierte
        Gesichter im Publikum suchen; nicht einfach das Programm abspulen.
  \item \textbf{Der theatralische Rednertyp.} Geht mit unvorhergesehenen
        Situationen spielerisch um, tritt locker auf und wirkt authentisch.
        Läuft Gefahr, das eigentliche Thema aus den Augen zu verlieren, und
        kann unsympathisch wirken. Emotionen in Maßen halten das Publikum
        wach und verhindern Monotonie, müssen aber zur Art der Rede passen.
  \item \textbf{Der distanzierte Rednertyp.} Zeichnet sich oft durch besondere
        Höflichkeit aus, läuft aber Gefahr, am Publikum vorbeizureden. Wirkt
        häufig selbstbewusst und kompetent, kann aber unnahbar wirken. Sollte
        die Distanz aufbrechen, ohne zu übertreiben: Nähe durch
        Verständnisfragen, Blickkontakt und Einbeziehung des Publikums.
  \item \textbf{Der überlegene Rednertyp.} Ausprägungen: der tatsächlich
        Überlegene (Siegertyp, in besonderem Maße begabt) und der unsichere
        Überlegene, der das Vortäuschen zum Selbstschutz nutzt. Versetzt das
        Publikum zurück in die Schulzeit und kann arrogant, überheblich oder
        gelangweilt wirken. Tipps: Sprechpausen für ein Lächeln ins Publikum
        nutzen, Blickkontakt in die Runde.
\end{itemize}

\subsection{Manipulation durch Rhetorik}

\begin{definitionbox}
\textbf{Rhetorisch manipulieren} heißt, ausgewählte Personen bewusst und
zielgerichtet unter Anwendung von Manipulationstechniken zu beeinflussen.
Ziel ist immer die \emph{verdeckte} Manipulation des Gegenübers; dadurch kann
das Verhalten einer Person willentlich beeinflusst werden.
\end{definitionbox}

Man sollte die grundlegenden Techniken kennen, um sich selbst davor schützen
zu können. Sie lassen sich auch im positiven Sinne einsetzen und geben
Aufschluss über Aspekte der menschlichen Psyche. Trotzdem gilt: vorsichtiger
und gewissenhafter Umgang.

\subsubsection{Techniken}

\begin{itemize}
  \item \textbf{Spiegeltechnik:} \glqq Gleich und Gleich gesellt sich
        gerne.\grqq{} Es findet ein Angleichen an Körpersprache, Mimik und
        Gestik des anderen statt. Spiegeln gaukelt Vertrauen und Sympathie
        vor; zeitverzögertes Spiegeln wird empfohlen.
  \item \textbf{Call to Action:} durch künstlichen Zeitdruck wird das
        Gegenüber in Bedrängnis gebracht, sich schnell entscheiden zu müssen.
        Oft werden nur Alternativen angeboten, die für das Gegenüber
        nachteilig sind.
  \item \textbf{Schlechtes Gewissen:} Mitleid ist die Anteilnahme am Unglück
        anderer und erhöht die Bereitwilligkeit zu helfen. Der Sprecher
        täuscht einen Gefühlsausbruch vor und stellt sich als hilfsbedürftig
        dar. Tricks: Kindchenschema, in eine Rolle schlüpfen,
        Emotionsworte verwenden.
  \item \textbf{Vermeidliche Experten:} der Sprecher beruft sich auf
        Wissenschaftler, Autoritätspersonen und namhafte Experten. Ob es
        diese Experten wirklich gibt, spielt dabei keine Rolle.
  \item \textbf{Assoziationstechnik:} durch Assoziationen wird beim Gegenüber
        ein bestimmter Emotionszustand hervorgerufen und es damit in eine
        Richtung gelenkt (Manipulation durch Visionen).
  \item \textbf{Sprachlicher Kontrast:} Kontraste entstehen durch
        Gegenüberstellung. Um herauszustechen, braucht man einen Vergleich,
        neben dem man gut dasteht.
  \item \textbf{Polarisierung:} dem anderen einen Grund liefern, etwas für
        einen zu tun, diesen Grund aber nicht verbalisieren, sondern nur in
        Aussicht stellen.
\end{itemize}

\subsubsection{Die dreizehn Scheinargumente}

\begin{center}
\begin{tabular}{ll}
Traditionsargument & Kompositionsargument\\
Autoritätsargument & Innovationsargument\\
Mehrheitsargument & Zirkelschlussargument\\
Entweder-Oder-Argument & Argument der penetranten Wiederholung\\
Argument auf die Person & Strohmann-Argument\\
Argument der Autorität & Argument des Vergleichs\\
Argument der Ignoranz & \\
\end{tabular}
\end{center}

Das Traditionsargument in Reinform: \glqq Das haben wir schon immer so
gemacht!\grqq{}

\begin{notebox}
\textbf{Take home message.} Machen Sie sich die Möglichkeiten Ihres Körpers
und vor allem Ihrer Stimme klar und setzen Sie diese bewusst ein. Deuten und
korrigieren Sie Körpersprache bei anderen und bei sich. Überzeugend wirkt vor
allem die Kongruenz zwischen Körpersprache und Gesagtem. Es werden sechs
kulturunabhängige Basisemotionen unterschieden. Und es gibt zahlreiche
rhetorische Tricks zur Manipulation: wer sie kennt, kann sich schützen.
\end{notebox}

% =========================================
\section{Storytelling und Elevator Pitch}
% =========================================

\subsection{Der Elevator Pitch}

\begin{definitionbox}
Der \textbf{Elevator Pitch} (\glqq Aufzugspräsentation\grqq{}) ist die
einprägsame Darstellung der Kurzfassung eines Businessplans, einer
Geschäftsidee oder der eigenen Stärken. Er wird meist ohne Hilfsmittel und in
freier Rede vorgetragen.
\end{definitionbox}

Wichtig ist das \textbf{Was}, nicht das \textbf{Wie}. Er fordert eine Menge
Konzentration und Übung, und die Begeisterung muss ansteckend sein.

\textbf{Zwei Längen:} der 30-Sekunden-Pitch für Networking, Events,
Vorstellungsrunden und Telefonate; der 3-Minuten-Pitch für Präsentationen vor
Venture Capitalists, Konferenzpräsentationen, Einstellungsgespräche und
Verkaufsmeetings.

\textbf{Herkunft:} in den 80er Jahren mussten Vertriebler mit ihren Chefs
Ideen und Lösungen innerhalb einer Liftfahrt besprechen; in den 90ern mussten
Unternehmensgründer die Entscheider der Venture-Capital-Gesellschaften
treffen.

\subsubsection{Die AIDA-Formel}

\begin{itemize}
  \item \textbf{A} Attention: Aufmerksamkeit erregen, Zielpersonen kurz
        ansprechen und dabei auffallen, aus der breiten Masse abheben.
  \item \textbf{I} Interest: Neugierde wecken, durch Bilder oder Beispiele
        erklären, neue Gesprächsstrategie, Spaß erzeugen, Konzentration auf
        das Wesentliche.
  \item \textbf{D} Desire: den Wunsch erwecken, auch etwas vom Kuchen
        abzubekommen, am Nutzen teilhaben zu können, den Herzschlag des
        Gesprächspartners erhöhen.
  \item \textbf{A} Action: Austausch der Visitenkarten mit der Vereinbarung
        eines weiteren Gesprächs, vielleicht schon eine Terminvereinbarung.
\end{itemize}

\subsubsection{Vorbereitung}

\textbf{Zielgruppe und Zielperson:} definieren und Inhalte und Inszenierung
darauf anpassen; der Pitch soll auf die Zielperson genau zugeschnitten sein,
also vorher Informationen einholen. Lässt sich keine konkrete Person nennen,
nimmt man eine fiktive.

\textbf{Ziel:} \glqq Was will ich erreichen?\grqq{} So konkret wie möglich
überlegen, was man von dieser bestimmten Person möchte, was erreichbar ist,
und das Ziel aufschreiben.

\textbf{Was hat das Gegenüber davon?} Warum soll er mir eine Chance geben?
Was kann mein Produkt für meinen Kunden leisten? Was sind die Nutzenvorteile?
Oft werden Gesprächspartner mit technischen Details gelangweilt oder
verwirrt. Die inhaltliche Ebene sollte auch die eigene Großmutter verstehen
können.

\textbf{Was unterscheidet uns?} Warum ist meine Idee besser als die des
Mitbewerbers? In welchen Bereichen bin ich der Erste, der Einzige, der
Schnellste, der Innovativste, der Kreativste? Dann kritisch prüfen: Ist das
Unterscheidungsmerkmal wirklich einzigartig, oder sieht der Kunde es als
selbstverständlich an? Was ist Standard?

\textbf{Wer ist unser Markt?} Den Markt so genau wie möglich einschätzen und
kleinere Teilmärkte auswählen, die am schnellsten Erfolg versprechen.

\subsubsection{The Pitch Canvas}

\begin{center}
\begin{tabular}{p{0.46\textwidth}p{0.46\textwidth}}
\multicolumn{2}{c}{\textbf{Simple Statement:} Sag in einem Satz, was du
machst}\\\hline
\textbf{Pain ($+$ Gain):} Erläutere das Problem deiner Zielgruppe &
\textbf{Product:} Erkläre deine Lösung\\
\textbf{Product Demo:} Zeige dein Produkt &
\textbf{What's Unique:} Was macht deine Idee einzigartig?\\
\textbf{Customer Traction:} Bringst du bereits deine PS auf die Straße? &
\textbf{Business Model:} Was ist dein Geschäftsmodell?\\
\textbf{Investment:} Wieviel Kapital benötigst du, und wofür? &
\textbf{Team:} Investoren investieren in Persönlichkeiten\\\hline
\multicolumn{2}{c}{\textbf{Call to Action:} Wie geht es weiter?}\\\hline
\multicolumn{2}{c}{\textbf{Und warum gerade du?}}\\
\end{tabular}
\end{center}

\subsubsection{Mögliche Inhalte}

Zielpersonen; Nutzenargumente; Unterscheidungsmerkmale; Markteinschätzung;
Idee, Produkt, Dienstleistung; Personen, Team, Bewerber; Ziel der
Aufforderung.

\begin{notebox}
Der Köder muss dem Fisch schmecken und nicht uns selber.
\end{notebox}

\subsubsection{Wie Menschen entscheiden}

Es gibt zwei Möglichkeiten, wie Menschen entscheiden: rational und emotional.
Das Verhältnis ist das eines Eisbergs: \textbf{1/7 rational}, \textbf{6/7
emotional}. Über der Wasseroberfläche liegen Fakten, Daten und Zahlen,
darunter die vier Grundbedürfnisse.

\begin{center}
\begin{tabular}{ll}
\textbf{Grundbedürfnis} & \textbf{Inhalt}\\\hline
Pride & Anerkennung, Bewunderung, Stolz\\
Profit & Geld, Wohlstand, Gewinn\\
Pleasure & Freude, Spaß, Genuss\\
Peace & Sicherheit, Zufriedenheit, Ruhe\\
\end{tabular}
\end{center}

Der Herzschlag ist die emotionale Bindung an eine Information. Er erhöht
sich, wenn der Zuhörer aus Tatsachen eine Verbindung zu sich selbst und die
damit verbundene Wirkung herstellen kann.

\subsubsection{Grundsätze bildhaften Verkaufens}

Zwei Formulierungen derselben Aussage bleiben unterschiedlich gut haften:
\glqq Die Implementierung von audio-visuellen Medien bei
Qualifizierungsmaßnahmen ist temporär umgekehrt proportional dem Einsatz von
multimedialen Projektionssystemen\grqq{} gegen \glqq Kennen Sie auch den
guten alten Overheadprojektor? Heute wird er kaum noch verwendet, weil wir
mehr bewegte Bilder und Geschichten mit einem Beamer zeigen können.\grqq{}

\begin{center}
\begin{tabular}{p{0.46\textwidth}p{0.46\textwidth}}
\textbf{Linke Gehirnhälfte} & \textbf{Rechte Gehirnhälfte}\\\hline
digitales, rationales, analytisches und logisches Denken; Regeln und Gesetze;
Wissenschaft; Lesen, Schreiben und Rechnen; linear; Zeit
&
analog; Emotionen, Kreativität, Musisches, Räumliches und Bilder;
ganzheitliches Problembewusstsein; Intuition; non-linear; Raum\\
\end{tabular}
\end{center}

Die Kunst ist es, beide Hälften zusammenarbeiten zu lassen. 80\,\% unseres
Wissens stammen aus Informationen, die wir über das Auge aufgenommen haben.
Im klassischen Elevator Pitch fehlt die Möglichkeit der Visualisierung, also
gilt: Kopfinformationen mit dem Bauch (den 4 Ps) verknüpfen. Jeder rationale
Nutzen soll seine bildhafte Entsprechung im Bauch finden; der
Gesprächspartner soll innerlich sehen, hören, fühlen, schmecken und riechen.

\subsubsection{Inszenierung}

Ein wirkungsvoller Elevator Pitch ist wie ein kleines Theaterstück, und ein
Theaterstück wirkt nur, wenn die gesamte Rhetorik stimmt. Darum wird es lange
vorher geprobt. Genau das gilt auch für den Pitch: inszenieren und proben,
und einen Spannungsbogen bauen.

\subsubsection{Die zehn goldenen Regeln}

\begin{enumerate}
  \item Seien Sie zuhörerorientiert.
  \item Beginnen Sie mit einem Interessen-Katalysator.
  \item Was ist der Grund, sich an Sie zu erinnern?
  \item Beschreiben Sie \emph{nicht} die Idee, das Produkt oder die
        Dienstleistung.
  \item Verwenden Sie eine bildhafte Sprache.
  \item Was ist die Problemlösung?
  \item Was hat der Andere davon?
  \item Das Ende: die Aufforderung zur Tat.
  \item Zeigen Sie Begeisterung.
  \item Übung macht den Meister.
\end{enumerate}

\subsection{Storytelling}

Storytelling heißt, Informationen in Form von Geschichten zu erzählen. Es
macht wichtige Informationen besser verständlich, unterstützt das Lernen und
Mitdenken der Beteiligten nachhaltig und fördert die geistige Beteiligung.

\subsubsection{Vier Aufgaben in den Public Relations}

\begin{enumerate}
  \item Es macht auf das Unternehmen aufmerksam.
  \item Es informiert über das Unternehmen.
  \item Es löst bedeutende Gefühle in den internen und externen
        Bezugsgruppen aus.
  \item Es sorgt dafür, dass die Bezugsgruppen das Unternehmen besser
        speichern und aus ihrem Gedächtnis leichter und schneller abrufen
        können.
\end{enumerate}

\subsubsection{Effekte von Storytelling und Storylistening}

\begin{center}
\begin{tabular}{ll}
\textbf{Struktur} & \textbf{Empathie}\\\hline
Fokus & Vertrauen\\
Motivation & Bindung\\
Erinnerung & Großzügigkeit\\
\end{tabular}
\end{center}

\subsubsection{Story-Struktur: B-T-E}

\begin{itemize}
  \item \textbf{B} Beginn: \glqq So wie es jetzt ist.\grqq{}
  \item \textbf{T} Transformation: \glqq Es passiert etwas.\grqq{}
  \item \textbf{E} Ende: \glqq Es hat Veränderung stattgefunden.\grqq{}
\end{itemize}

\subsubsection{Wandel anregen}

Probleme ansprechen und zeigen, dass es Gewinner und Verlierer geben wird.
Die meisten Menschen neigen zur \textbf{Verlustaversion}: Verlust wird höher
gewichtet als Gewinn. Sie wird überwunden, indem man demonstriert, wie der
Wandel große Gewinner und große Verlierer hervorbringt: der Wandel in der
Story wird so angepasst, dass er wahrscheinlich in einer sehr positiven
Zukunft für das Gegenüber endet, während Nichtstun in einer inakzeptabel
negativen Zukunft endet.

% =========================================
\section{Wissenschaftliches Arbeiten}
% =========================================

\subsection{Der wissenschaftliche Blick}

\begin{definitionbox}
\textbf{Vier Prinzipien wissenschaftlichen Arbeitens:} Ehrlichkeit,
Objektivität, Überprüfbarkeit, Nachvollziehbarkeit.
\end{definitionbox}

\subsection{Literaturarbeit}

Die Suche läuft über Schlagworte und boolesche Operatoren, und die Wahl des
Operators entscheidet über die Größenordnung des Ergebnisses.

\begin{examplebox}
Suche nach \emph{University}, \emph{Industry}, \emph{Collaboration}:
\begin{center}
\begin{tabular}{lr}
\textbf{Verknüpfung} & \textbf{Treffer}\\\hline
mit OR & 888.385\\
mit AND & 1.236\\
als exakte Phrase & 211\\
\end{tabular}
\end{center}
\end{examplebox}

Zur Beurteilung von Quellen dienen \textbf{Journal Rankings} (Impact Factor,
Financial Times 50, VHB-Ranking) und die \textbf{Zitationshäufigkeit}. Wie
oft ein Artikel zitiert wurde, kann Aufschluss über seine Relevanz geben;
Achtung: neuere Artikel sind natürlich weniger zitiert.

\subsubsection{Zitierwürdig und zitierfähig}

\begin{center}
\begin{tabular}{p{0.3\textwidth}p{0.3\textwidth}p{0.3\textwidth}}
\textbf{zitierwürdig} & \textbf{nicht zitierwürdig} &
\textbf{z.T.\ zitierwürdig, beschränkt zitierfähig}\\\hline
Originalarbeiten;
wissenschaftliche Fachbücher;
Dissertationen;
Artikel in Fachzeitschriften;
(Electronic) Journals
&
nicht wissenschaftliche Bücher;
allgemeine Lexika;
Einführungsliteratur;
Skripte;
Seminararbeiten;
Artikel in Boulevardzeitungen;
allgemeine Internetseiten (inkl.\ Wikipedia)
&
Firmenschriften;
Diplomarbeiten;
graue Literatur;
Fachwörterbücher und Fachlexika\\
\end{tabular}
\end{center}

\subsection{Zitation}

Richtiges Zitieren ist das Handwerkszeug des Wissenschaftlers. Der Autor macht
deutlich, welche Gedanken und Befunde von anderen übernommen wurden. Zweck
ist immer die Nachprüfbarkeit und Vertiefung für andere. Plagiatsversuche
können bis zur Exmatrikulation führen.

\begin{center}
\begin{tabular}{p{0.46\textwidth}p{0.46\textwidth}}
\textbf{Direktes Zitat} & \textbf{Indirektes Zitat}\\
Quelle wird wörtlich angeführt & Quelle wird sinngemäß angeführt\\\hline
\ldots fasst Owens (1976, S.~625) wie folgt zusammen: \glqq Unser Grundaxiom
lautet, dass der beste Prädiktor künftigen Verhaltens in der Vergangenheit
gezeigtes Verhalten ist.\grqq{}
&
Zukünftiges Verhalten lässt sich am besten durch in der Vergangenheit
gezeigtes Verhalten prognostizieren (vgl.\ Owens 1976, S.~625).\\
\end{tabular}
\end{center}

In den Wirtschaftswissenschaften sind indirekte, sinngemäße Zitate die Regel.
Ein umformuliertes Zitat, das dem Original zu nahe bleibt, ist trotz
Quellenangabe ein Plagiat.

\begin{examplebox}
\textbf{Original:} \glqq Bei der einfachen Zufallsstichprobe hat jedes
Element der Grundgesamtheit die gleiche Chance, ausgewählt zu werden
[\ldots].\grqq{}

\textbf{Plagiat:} \glqq Bei der einfachen Zufallsauswahl haben alle Elemente
der Grundgesamtheit die gleiche Chance, ausgewählt zu werden [\ldots]\grqq{}
(vgl.\ Ebster 1999, S.~82).

\textbf{Korrektes indirektes Zitat:} \glqq Für alle Untersuchungseinheiten
besteht bei der einfachen Zufallsauswahl die gleiche Wahrscheinlichkeit, in
die Stichprobe aufgenommen zu werden\grqq{} (vgl.\ Ebster 1999, S.~82).

Es geht um die eigene intellektuelle Leistung.
\end{examplebox}

\subsubsection{Belegweise}

Es gibt zwei Möglichkeiten, Quellen im Text zu belegen: als Anmerkungen in
Fuß- oder Endnoten (im deutschen Raum verbreitet) oder als Kurzbelege im Text
(im angloamerikanischen Raum verbreitet), also \glqq Kundenzufriedenheit führt
zu Kundenloyalität (Oliver, 1999, S.~35).\grqq{}

\subsection{Formalitäten}

\textbf{Inhaltliche Gestaltung:} die eigene Meinung zurückhalten; auch
allgemeine Aussagen belegen; aktueller Bezug und Praxisorientierung;
eigenständige Darstellung.

\textbf{Sprachliche Gestaltung:} knapp und nüchtern, unaufgeregt und
sachlich; Fremdwörter korrekt benutzen; Wortzusammensetzungen vermeiden;
blumige Sprache vermeiden; gegenlesen lassen, Rechtschreibung und Grammatik
prüfen.

\textbf{Titel der Arbeit:} die Visitenkarte der Arbeit. Der Titel sollte den
Fokus klar beschreiben, auch wenn er dann weniger griffig klingt. Untertitel
können den Fokus weiter verdeutlichen.

\begin{examplebox}
\glqq Fortsetzung folgt\ldots{} Der Effekt von Cliffhangern bei narrativen
Serien auf Binge-Watching am Beispiel von Netflix\grqq{}
\end{examplebox}

Alle weiteren Formalitäten (Deckblatt, Inhaltsverzeichnis, Schriftart,
Zitationsstil) sollten mit der betreuenden Person abgesprochen werden.

\subsection{Zeitmanagement}

\textbf{Planung:} Ziele genau definieren; den Zeitrahmen realistisch
einschätzen (durch Selbstbeobachtung und Durchschnittswerte); mit ToDo-Listen
arbeiten (heute, in 2--4 Tagen, in 5--10 Tagen, später).

\subsubsection{Arbeitsschritte einer wissenschaftlichen Arbeit}

\begin{enumerate}
  \item Themensuche
  \item Grundrecherche
  \item Grobgliederung
  \item Literaturrecherche und -auswertung
  \item Forschungsdesign, Datenerhebung, Auswertung
  \item Texterstellung
  \item Fertigstellung: Layout, Druck, Kurzfassung (Abstract)
\end{enumerate}

Zur Terminierung dient ein \textbf{Gantt-Diagramm}: die Aktivitäten stehen in
den Zeilen, die Kalenderwochen in den Spalten, und ein Balken markiert, wann
welche Aktivität läuft.

\subsubsection{Motivation durch Belohnung}

Grundlage der modernen Verhaltenstherapie, operante Konditionierung. Zwei
Möglichkeiten, die parallel genutzt werden sollten: das Arbeitsumfeld positiv
gestalten und die geleistete Arbeit belohnen (Urlaub, Entspannung, Sport).

\subsubsection{Eisenhower-Prinzip}

\begin{center}
\begin{tabular}{lll}
 & \textbf{dringend} & \textbf{nicht dringend}\\\hline
\textbf{wichtig} & sofort selbst erledigen & terminieren\\
\textbf{nicht wichtig} & delegieren & Papierkorb\\
\end{tabular}
\end{center}

% =========================================
\section{Zur Klausur}
% =========================================

\begin{itemize}
  \item Format: Single Choice, 30 Fragen zu je einem Punkt, 60 Minuten.
  \item Keine Hilfsmittel zugelassen.
  \item Die Studienleistung ist \textbf{unbenotet}, sie muss nur bestanden
        werden.
  \item Klausurrelevant sind alle behandelten Themen der Vorlesung. Nicht
        klausurrelevant sind die Inhalte der Gastvorträge und die Folien zur
        Bedienung des Folienmasters (aus \glqq 4. Präsentation\grqq{}, ab
        Folie 51).
  \item Inhalte, die im Wrap-Up nicht erwähnt werden, können trotzdem
        Bestandteil der Klausur sein.
\end{itemize}

\subsection{Was am ehesten geprüft wird}

Die Zahlen und die benannten Listen sind das, was sich in einer Single-Choice-
Frage abfragen lässt:

\begin{itemize}
  \item die fünf Axiome von Watzlawick und die vier Seiten einer Nachricht;
  \item die sechs Ich-Zustände und die drei Transaktionsarten;
  \item die drei Elemente des Feedbacks;
  \item 41\,\% Angst vor öffentlichem Reden; die drei Stressmodelle mit ihren
        Autoren; die drei Phasen des Adaptionssyndroms;
  \item SMART; die Aufteilung 15/75/10; die vier Phasen der Inhaltsplanung;
  \item die Behaltensquoten 20/50/100; die vier Verständlichmacher; 16\,pt
        für 5\,m und 24\,pt für 10\,m; höchstens drei Farben; 5 bis 7 Punkte
        und höchstens sieben Folien am Stück;
  \item Ethos, Pathos, Logos; die vier Distanzen; die sechs Basisemotionen;
        die vier Pausenarten; die fünf Rednertypen; die dreizehn
        Scheinargumente;
  \item AIDA; 30 Sekunden gegen 3 Minuten; 1/7 gegen 6/7; die vier
        Grundbedürfnisse (Pride, Profit, Pleasure, Peace); die zehn goldenen
        Regeln; B-T-E;
  \item zitierwürdig gegen zitierfähig; direktes gegen indirektes Zitat; die
        sieben Arbeitsschritte; das Eisenhower-Prinzip.
\end{itemize}

\end{document}
